\begin{annexitem}{annex:p12.eb.ground-termination}{thm:p12-eb-ground-termination}{Every open rewrite step has a finite-support grounded instance, including stutters}
\begin{proof}
Choose an original ground term q and the termwise map H_q.
Assume an infinite open run \(t_0\to_Rt_1\to_R\cdots\).
For its i-th step write
\(t_i=C_i[l_i\gamma_i]\),
\(t_{i+1}=C_i[r_i\gamma_i]\), with the chosen rule variant.
There are finitely many variables in \(l_i,r_i\). Define the
finite substitution \(\mu_i\) on those rule variables by
\(x\mu_i=H_q(x\gamma_i)\), identity elsewhere. Homomorphic
evaluation gives
\[
 H_q(t_i)=H_q(C_i)[l_i\mu_i]
       \quad\longrightarrow_R\quad
 H_q(C_i)[r_i\mu_i]=H_q(t_{i+1}).
\]
Both terms and the surrounding context are ground. Every step
is a valid original-signature contextual instance. This works
even if the run introduces infinitely many different variables:
each \(\mu_i\) has finite support, and H_q is only a proof map.
If two neighboring H_q-images coincide, the same rule instance
is a permitted equal-endpoint rewrite step. Repeating the steps
therefore still produces an infinite ground derivation, contradicting
ground termination. Open termination follows.

Neither the existence nor the absence of such self-steps can be
silently changed. For example, over the constant-only signature
\(\{a\}\), the unrestricted rule x->y has the ground instance
a->a. Under the source definition its ground relation does not
terminate. If one instead discarded all equal-endpoint instances,
its ground relation would be empty while its open variable-to-variable
relation would have infinite runs. This is a different convention,
not a rule respecting a reduction ordering in Theorem4.13.
Similarly, without a ground term the implication is false: over
one unary f and no constants, f(x)->f(f(x)) has no ground terms
at all and an infinite open run. It too does not satisfy the
source's separate reduction-ordering premise.

The map H_q is not injective: H_q(x)=H_q(q)=q. It does not
reflect open equality or unifier coverage. The proof establishes
termination only and uses neither distinct parameter constants
nor an unfreezing operation.
\end{proof}
\end{annexitem}

\begin{annexitem}{annex:p12.eb.initial-normalization}{thm:p12-eb-initial-normalization}{The order supplies ordinary rule form and finite open-image normalization}
\begin{proof}
For a rule l->r, full invariance gives
\(l\gamma\succ r\gamma\) for every finite substitution;
monotonicity preserves this direction in every context. Every
rewrite step therefore decreases the well-founded order on the
whole open term domain, proving termination there. This argument
also covers signatures without any ground term.

Suppose a variable z occurs in r but not in l. Substitute l for z.
The left side stays l, while the right side either becomes l
itself (if r=z) or contains l inside a proper context D.
Strict order would then give either \(l\succ l\) or
\(l\succ D[l]\). In the latter case monotonicity iterates
the fixed context to give an infinite descending sequence
\(l\succ D[l]\succ D[D[l]]\succ\cdots\), impossible.
Thus every right-side variable occurs on the left.
If l is a variable x and r omits x, substituting r for x
gives a strict self-comparison. If r is x the rule already
violates strictness. If r properly contains x, the same context
iteration excludes \(x\succ r\). Hence l is nonvariable.
This derives the ordinary-rule form from the actual ordering
assumption rather than adding a hidden syntactic restriction.

For each of the finitely many nonidentity images of \(\theta\),
choose redexes until a normal form is reached; open termination
makes every such process finite. Put that form at the same
variable in \(\theta_0\), dropping an identity binding if it
arises. Outside the old finite support leave variables fixed.
Those variable terms are R-normal because no R-left-side is a
variable. Therefore \(\theta_0\) has finite support and is
R-reduced. Each chosen reduction gives
\(x\theta=_Rx\theta_0\); outside the support this equality
is identity. Induction on terms, using congruence, gives
\(s\theta=_Rs\theta_0\) for every finite s. Thus every
input equation solved by \(\theta\) is still solved by
\(\theta_0\). Both problem-variable instantiation comparisons
use the identity factor. This does not require the normalized
images to be mutually substituted or the target to be idempotent.

No claim of unique normal forms has been used. Ground confluence
does not choose unique normal images of open terms; choose any
finite image reductions. Since initial S is empty, this construction
imposes no changed literal binding. It is not normalization of a
reachable active accumulator.
\end{proof}
\end{annexitem}

\begin{annexitem}{annex:p12.eb.open-history-lift}{thm:p12-eb-open-history-lift}{Open-address compatibility, computed U factor and every protected origin}
\begin{proof}
The primitive solved-state invariant gives
\(\sigma_S\rho=\rho\) globally. A literal C elimination
x=t has \(x\rho=t\rho\), so its one-variable map
\(\delta\) satisfies \(\delta\rho=\rho\). Updating C,S
therefore preserves this same witness; B leaves P unchanged.
Trivial, Orient and Decomposition likewise preserve literal
solutions and change none of the witness images. The original
comparison scope remains external, so deletion of a variable's
last equation does not delete its target assignment or factor.

If a chosen forward redex in s\(\rho\) lay at or inside an
inserted image x\(\rho\), that image would be R-reducible.
Every image in the support is normal; every remaining variable
is also normal by the derived ordinary-rule form. Hence the
redex is at a nonvariable skeleton address p. Let u=s|_p,
choose the innermost redex and a fresh rule l->r with
\(l\gamma=u\rho\). Choose the variant outside the earlier
computation and the finite support of \(\rho\).
Every matched variable image is a proper subterm of this
innermost redex and is normal, even when the subterm contains
parameters. Repeated variables receive the same value because
this is an actual match. There are no unassigned right-only
variables in the forward rule.

Extend \(\rho\) by those fresh normal images, obtaining
\(\rho'\). Freshness preserves all old values and literal
C,S solutions. Its matching equation
\(Q=\{l\sigma_S=^{?}u\sigma_S\}\) is literally solved:
both sides under \(\rho'\) are l\(\gamma\)=u\(\rho\).
When l is nonvariable its root matches u; first Decomposition
is legal and is performed before any other operation on this
new equation. Run the actual removing-binding U helper on the
resulting literal constraints. Since \(\rho'\) solves them,
the finite run succeeds and its actual \(\eta\) satisfies
\(\eta\rho'=\rho'\) globally. Its domain/range are among
variables of Q and avoid old solved domains. The actual new
bindings are \(S\eta\cup[\eta]\), with substitution
\(\sigma_S\eta\), so
\(\sigma_S\eta\rho'=\rho'\). Existing duplicate literal
constraints are reconciled by Trivial as in the helper-block
interface; C-empty block boundaries agree exactly.

The new skeleton e[r] under \(\rho'\) is precisely the
chosen original contextual reduct, and all other old equation
values are unchanged. E-equality preserves their solution
conditions; literal helper equations and S are solved by the
same explicit factor. Because every actual R-step decreases
\(\succ\), the problem-value multiset decreases for this
chosen forward block. No groundness or confluence was used
to prove this local assertion.

Mark function occurrences of initial P as skeleton occurrences.
At LP, keep the unchanged context marks and mark all function
occurrences of the fresh right-side template. Constrain deletes
marks with its removed equation; C operations and Basic
Elimination change none of P's marks. Induction through the
actual six rules proves that these are exactly the available
nonvariable skeleton addresses. Any functions appearing only
when a variable is instantiated are below that variable's
address and have no skeleton mark. Under the maintained witness
their whole enclosing image is normal; thus no selected forward
redex can be one of those protected functions. This is an
actual origin invariant, not an assumed successful basic proof.

Also, each solved right side t has
\(t\rho'=x\rho'\), the normal image of its solved left
variable. If t had an R-redex, rewriting's substitution stability
would give a redex in that normal image. This remains true for
an equal-endpoint instance because the source counts it as a
step. Thus the computed accumulator itself has reduced images
along this compatible witness branch; it has not been normalized
in place. An inverse expansion can still match a forward-normal
image. Its exclusion from protected positions comes from the
origin invariant, not from forward normality. Nothing in this
proof ensures that an unsolved E-equal pair has a forward redex.
\end{proof}
\end{annexitem}

\begin{annexitem}{annex:p12.eb.inverse-expansion}{thm:p12-eb-inverse-expansion}{Normal inverse matching, fresh right-only assignments and the actual increase}
\begin{proof}
Use the specified eligible u and fresh reversed r=l. Its
match r\(\gamma\)=u\(\rho\) is given, not assumed to
exist universally. If r is nonvariable, each variable appearing
in it receives a proper subterm of the normal u\(\rho\),
hence a normal term. If r is a variable, its single image is
the whole normal u\(\rho\). Choose every remaining variable
of l a fresh parameter variable, which is normal for ordinary R.
This extends \(\gamma\) to all variables of the identity with
normal images. Choose all variant and parameter names outside
the finite old computation/support. The resulting \(\rho'\)
preserves old values and literally solves the old C,S.

Its new literal constraint is
\(r\sigma_S=^{?}u\sigma_S\); both sides evaluate to the
given match. For nonvariable r, roots match and first matching
Decomposition is performed. For variable r that obligation is
absent. The actual U helper succeeds and yields
\(\eta\rho'=\rho'\) and new solved substitution
\(\sigma_S\eta\), with the same global factor calculation
as above. The identity r=l ensures E-solution of the new
problem equation while all old variable values stay fixed.

The new nonvariable template is l; its function occurrences
are therefore genuinely introduced skeleton marks. Substitution
images of all its fresh variables remain normal, so any forward
redex created by expansion is at an eligible marked address.
Finally, the forward rule has \(l\succ r\). Full invariance
gives \(l\gamma\succ r\gamma=u\rho\), and monotonicity
gives the same strict increase in the enclosing problem term.
Thus this inverse block does not decrease the source multiset M.
This proves its compatibility and history assertions without
claiming a global progress measure or a complete inverse strategy.
\end{proof}
\end{annexitem}

\begin{annexitem}{annex:p12.eb.progress-countercase}{thm:p12-eb-progress-countercase}{Original ground convergence, fixed normal identity witness and a complete example branch}
\begin{proof}
Every term is a unary f/g word above either a or one variable.
Let \(\operatorname{len}(t)\) count all its nodes, including
the terminal leaf. Define t\(\succ\)u when
\(\operatorname{Var}(u)\subseteq\operatorname{Var}(t)\)
and either t has greater length, or lengths are equal and the
f/g prefix word of t is lexicographically greater with f>g.
Variable/constant terminal leaves are not ordered by this tie test.
This is a strict partial order: variable inclusion is transitive,
length comparisons are transitive, and equal-length prefix words
use a transitive lexicographic order. Each descendant has length
no greater than its predecessor and may use only its predecessor's
variable. Below an initial finite length there are only finitely
many such words and terminal choices. Thus no infinite strict
descent exists. On ground terms the leaf is always a; different
equal-length terms have different prefix words. Hence the order
is total on ground terms.

Contexts are common unary prefixes, so both length and the first
different prefix symbol are preserved in a context. For substitution
stability, if both terms have a variable, inclusion forces the same
one; replacing it adds the same number of nodes to both terms and
changes nothing before their first prefix difference. If only t
has a variable, u is ground. Replacing that leaf adds a nonnegative
number of nodes to t; a previous length advantage stays strict.
A tied prefix advantage either stays at the same first different
function symbol, or becomes a strict length advantage. The case
of a variable only on the right is excluded by inclusion. Ground
terms are unaffected. This proves full invariance and monotonicity.
The first rule is oriented by the tied prefix comparison f>g;
the other two rules strictly shorten the word. Thus R satisfies
even a ground-total reduction-ordering reading of the source premise.
All its left sides have at most one occurrence of each variable,
and each rule has the same variable set on both sides.

Every closed term over this signature is a chain of f,g above a.
Normalize its child by induction on its tree. An f-parent rewrites
by the first rule to g of that normalized child; g(a) rewrites
to a. Hence each closed term has a reduction to a. Moreover
an irreducible closed term cannot be rooted at f, and a closed
g-term with irreducible child a is also reducible. So a is the
only closed normal form. Termination then proves ground
confluence in the original signature.

The open peak f(g(y)) contracts at its root by the first rule
to g(g(y)), or by the second to g(y). Those terms are distinct
and irreducible: neither contains f or a. They are nevertheless
R-equal by this peak, so identity is the specified reduced
unifier. The source proof's C-case is unavailable, literal
equality is unavailable, and no forward redex exists.

Any R-normal term over this signature is a, a variable z, or
g^k(z) for some k>=1. A two-element model with a=0 and f=g
constantly0 satisfies all three identities. Give y value1.
It distinguishes y from a and every g^k(z). If z is a
different variable, choose its value0 and y's value1 in that
same model. Thus the only R-normal term R-equal to y is y.
Every alternative reduced initial witness equivalent to identity
on y therefore has exactly the same blocked endpoints. The
obstruction is not remedied by a different image-normalization choice.

For the actual raw inverse-basic branch, use g(v)=f(v) at
the left g(g(y)) root. First matching Decomposition gives
v=g(y); Basic Elimination records S={v=g(y)} and leaves
the replacement skeleton f(v)=g(y). Extend the witness by
v\(\rho\)=g(y), a normal open image. Next use
f(g(w))->g(w) at the f-root. The match after \(\sigma_S\)
is f(g(w))=f(g(y)). Its mandatory first Decomposition gives
g(w)=g(y); the next literal decomposition and elimination
give w=y. The new skeleton is g(w)=g(y), and
S={v=g(y),w=y}. Constrain instantiates it to g(y)=g(y);
Trivial ends the trace. Both new witness images are normal,
the original y is fixed, every LP address is a genuine skeleton
function occurrence, and no active binding was normalized.

The returned substitution fixes y. Since the input equation is
an R-identity, every substitution solves it; this returned
singleton therefore covers every target on the initial scope
{y}, using that target as the factor. The example proves no
general inverse-basic theorem, but its branch is fully complete
for this particular problem.

At the first inverse step the problem-value multiset changes
from {g(g(y)),g(y)} to {f(g(y)),g(y)}. The first rule
implies f(g(y))\(\succ\)g(g(y)); hence that step increases
M and cannot be justified by source Theorem4.13's stated
descending multiset induction. Finally, adding a fresh constant c
would leave both g(g(c)) and g(c) irreducible, exposing the
same unjoined peak at f(g(c)). Ground confluence of the
original signature therefore does not silently extend to that
grounding vocabulary. Grounding y merely to the old a hides
the peak; replacing that existing rule constant back by y
would change g(a)=a, not preserve the original theory.
\end{proof}
\end{annexitem}
